% !TEX program = pdflatex
%
% FRONT-END DRAFT. This is a starting skeleton, not a submission-ready
% manuscript: swap the document class for the target journal's template
% (e.g. elsarticle for Journal of Transport Geography / Transport
% Reviews), move \thebibliography into a proper .bib file, and fill in
% author/affiliation details. Content is drawn from the development
% history of the IKOB Aporometrica accessibility-modelling codebase
% (https://github.com/Marrekoo/IKOB_Aporometrica) and should be checked
% against that history before submission -- in particular, the three
% incidents in Section 3 should be verified against the actual commit
% log / PR history rather than taken on faith from this draft.

\documentclass[11pt]{article}

\usepackage[margin=1in]{geometry}
\usepackage{amsmath,amssymb}
\usepackage{natbib}
\usepackage{hyperref}
\usepackage{booktabs}
\usepackage{enumitem}
\usepackage[T1]{fontenc}
\usepackage[utf8]{inputenc}

\title{Beyond Twelve Tips: What an Agentic Coding Assistant Catches in Applied Transport Modelling}

\author{
  Marco van Burgsteden\thanks{Corresponding author.} \\
  \small Department / Institute, placeholder \\
  \small \texttt{email@placeholder}
}

\date{}

\begin{document}
\maketitle

\begin{abstract}
\noindent
Recent guidance on AI-assisted scientific coding -- most notably Bridgeford
et al.'s twelve tips for AI-assisted coding in science -- treats the
large language model (LLM) as a coding instrument: it writes functions,
it writes tests, and the researcher remains responsible for correctness.
This framing fits code-centric science well, but it understates what
changes when an \emph{agentic} coding assistant also manages the
development workflow itself -- running tests, driving version control,
and flagging problems unprompted -- in a discipline where correctness is
not only a property of code but a property of \emph{fidelity to a
published specification}. We report on a short case drawn from
reconstructing a legacy transport-accessibility model (IKOB) as a
tested, documented Python package under a specification-first,
test-driven, agentic workflow. We distinguish three classes of issue an
agentic assistant surfaced during this work -- correctness-class,
governance-class, and specification-fidelity-class -- and argue that
only the first two are addressed by existing AI-assisted-coding
guidance. The third, where an agent flags that the \emph{implemented
model family} silently diverges from the \emph{specified} one, is
specific to applied modelling disciplines and is, to our knowledge, not
yet reported. We propose a minimal addition to existing AI-coding
checklists for modelling research: an explicit, agent- or human-checked
correspondence table between paper equations and the code objects that
implement them.
\end{abstract}

\noindent\textbf{Keywords:} AI-assisted coding; agentic LLMs; transport accessibility modelling; reproducibility; research software

\section{Introduction}
\label{sec:intro}

Guidance on using large language models (LLMs) for scientific software
has matured quickly. \citet{bridgeford2026} distil twelve practical tips
-- organised around problem preparation, context management, testing,
and code quality -- for using AI coding tools without sacrificing
scientific rigor. Their framing, like most of the surrounding literature
\citep{sandve2013,wilson2014,poldrack2024}, treats the LLM as an
instrument for \emph{writing and reviewing code}: a tool whose output
the researcher must independently verify against tests, domain
knowledge, and critical reading.

This framing is incomplete for a workflow that is increasingly common:
an \emph{agentic} coding assistant (e.g.\ Claude Code, Codex) that does
not just write functions on request but also runs the test suite,
drives git and continuous-integration workflows, monitors its own
output, and -- in our experience -- proactively flags things the
researcher did not ask it to check. \citet{bridgeford2026} acknowledge
this shift (their Tip~9, ``monitor progress and know when to restart'',
and Section~3.3, ``guardrails for autonomous agents''), but their
treatment stays at the level of code correctness and workflow safety.
Applied modelling disciplines such as transport geography have a third
failure mode that generic coding guidance does not cover: a piece of
code can be internally correct, fully covered by passing tests, and
still implement the \emph{wrong model}, because unit tests check
behaviour, not correspondence to a published equation.

We report a short case from rebuilding a legacy Dutch accessibility
model (IKOB) as a specified, tested, documented Python package, under a
protocol combining (i) a proven specification and codebase as the
starting point, (ii) specification-first, test-driven development, and
(iii) an agentic coding assistant used throughout, including for run
management and version control -- not only for writing functions. Over
the course of this work, the assistant surfaced issues in three
distinct classes, summarised in Table~\ref{tab:typology}. We argue that
the third class is the one with no existing home in the AI-assisted
coding literature, and that it matters specifically \emph{because} the
discipline's object of concern is not the code but the model the code
is supposed to represent.

\section{Setting}
\label{sec:setting}

The codebase under development reconstructs a competition-adjusted
gravity accessibility model (Hansen/Shen-style) with population
segmentation by income and household type, and a joint
time--money tolerance filter composed via a survival copula --
the implementation vehicle for a companion methodological paper
\citep{vanburgsteden-thresholds}. Development proceeded under an
explicit protocol: begin from a stated mathematical specification and a
previously working legacy codebase; rewrite incrementally with an
agentic LLM assistant, pairing every behavioural change with a unit
test before or alongside the change; keep a continuous, version-controlled
record of what changed and why (commit messages, pull requests, and an
explicit decision log). The assistant was given broad but
directory-scoped permissions: it could read and edit the repository,
run the test suite, and manage git operations, but version-control
history was always inspected before any destructive operation, in line
with \citeauthor{bridgeford2026}'s guardrail recommendations
(\S\ref{sec:intro}; \citealp{bridgeford2026}, \S3.3).

\section{A typology of catches}
\label{sec:typology}

\begin{table}[t]
\centering
\caption{Three classes of issue surfaced by the agentic assistant during development, and their treatment in existing AI-assisted-coding guidance.}
\label{tab:typology}
\small
\begin{tabular}{@{}p{2.6cm}p{5.6cm}p{5.4cm}@{}}
\toprule
Class & What happened & Covered by existing guidance? \\
\midrule
Correctness &
A numeric-column auto-detection routine crashed on real-world input
because a pandas/GDAL nullable dtype is not a plain \texttt{numpy}
dtype. A regression test reproducing the exact failure was written
before the fix, confirmed to fail on the old code, and confirmed to
pass after. &
Yes -- directly matches Tip~7 (\citeauthor{bridgeford2026}): use a
failing test to specify and guard against a known bug. \\
\addlinespace
Governance &
A merged pull request's branch diverged from \texttt{main} after a
local IDE sync, triggering a stop-hook warning about unpushed commits.
Before force-pushing to reconcile, branch and \texttt{main} contents
were diffed to confirm no work would be discarded, and the GitHub API
was queried to confirm the PR had in fact merged rather than vanished. &
Yes -- matches the guardrails in \S3.3 of \citeauthor{bridgeford2026}
(commit before changes, inspect history before anything destructive),
and standard version-control practice predating AI coding entirely. \\
\addlinespace
Specification fidelity &
While answering an unrelated question, the assistant noted that the
dependence structure implemented in the copula-composition module
(a Frank copula) did not match the one specified in the governing
methodological paper (a Gumbel--Hougaard survival copula). Both are
valid copula families and the code's unit tests -- which check the
Fr\'echet--Hoeffding bounds and other family-agnostic invariants --
passed regardless of which family was implemented. &
No -- not addressed by \citeauthor{bridgeford2026} or related
guidance, which treat correctness as a property of code behaviour
rather than of correspondence to a specific published equation. \\
\bottomrule
\end{tabular}
\end{table}

The first two classes in Table~\ref{tab:typology} are reassuring in a
specific sense: they show that an agentic assistant, used under a
test-driven and version-control-disciplined protocol, behaves
essentially as \citet{bridgeford2026} describe -- useful, occasionally
risky, and manageable with ordinary software-engineering guardrails.
The third class is different in kind, not degree. A copula family is a
modelling choice with first-order consequences for a paper's
conclusions (in our case, the policy interpretation of a dependence
parameter $\theta$ across income segments), yet nothing in a
conventional software test suite is positioned to catch its silent
substitution: the code was not \emph{broken}, it was \emph{unfaithful},
and unfaithfulness of this kind produces no stack trace.

\section{Discussion}
\label{sec:discussion}

Why did this catch happen at all? Not because the assistant was
specifically instructed to audit model fidelity, but because it had
simultaneous access to both the governing paper and the codebase across
an extended working session, and treated a mismatch between the two as
worth surfacing even unprompted -- a behaviour closer to a junior
collaborator cross-checking a derivation than to a code-completion
tool. This is precisely the capability that distinguishes an agentic
assistant used across an entire modelling workflow from the
code-generation-on-request pattern that \citet{bridgeford2026} and most
of the surrounding literature implicitly assume. It is also, we think,
why the guidance in that literature under-covers the risk: a
specification-fidelity failure is invisible to any checklist item
framed in terms of tests, documentation, or code review, because all
three operate \emph{within} the code's own behavioural contract, not
against the paper the code is meant to implement.

We do not think this generalises to a claim that agentic assistants
reliably catch specification drift -- our evidence is a single case,
surfaced opportunistically rather than systematically, and absence of
further catches in this project is not evidence of their absence
elsewhere in the code. The actionable claim is narrower: applied
modelling disciplines should not assume that the existing AI-assisted
coding checklists (document the workflow, test the code, review the
diff) are sufficient, because none of their items are built to catch
this failure mode, and a modelling paper's validity rests on it more
than most scientific software does.

\subsection{A minimal addition to existing guidance}

We propose a small, low-cost extension to \citeauthor{bridgeford2026}'s
Tip~12 (documentation for reuse and publication), specific to
quantitative modelling papers: alongside the existing checklist items
(workflow documentation, docstrings, environment specification, README,
worked examples), maintain an explicit \emph{equation--code
correspondence table} mapping each numbered equation in the
methodological paper to the function, class, or parameter that
implements it, and have either a human co-author or the coding
assistant itself check this table against the current state of the
code as a discrete, named step -- not folded silently into general code
review. This does not require new tooling; it requires only that
fidelity-checking be made a visible, repeatable step rather than an
incidental by-product of a long working session, as it was in our case.

\section{Conclusion}
\label{sec:conclusion}

Existing guidance for AI-assisted scientific coding treats correctness
as a property of code. In applied transport and geographic modelling,
correctness is also a property of a model's fidelity to its own
specification, and this is a failure mode that current AI-coding
checklists do not address. An agentic assistant working across an
entire modelling workflow -- not just generating functions on request
-- is positioned to catch this failure mode in a way that purely
code-centric tooling is not, but only if researchers treat
specification fidelity, not just test coverage, as something to
actively check for. We offer the typology in Table~\ref{tab:typology}
and the correspondence-table proposal in \S\ref{sec:discussion} as a
starting point for a more systematic treatment.

\subsection*{AI usage disclosure}
An agentic LLM coding assistant (Claude Code, Anthropic) was used
throughout the development of the codebase discussed in this note,
including for code authorship, test authorship, version-control
operations, and -- as discussed in \S\ref{sec:typology} -- for
surfacing the specification-fidelity issue reported here. The
assistant was also used in drafting this manuscript. All code, tests,
and commit history are available at the project repository
(placeholder: insert public URL at submission time, if the repository
is made public; it is currently private).

\subsection*{Data and code availability}
Code: \url{https://github.com/Marrekoo/IKOB_Aporometrica} (currently
private; visibility to be confirmed before submission).

\bibliographystyle{plainnat}
\begin{thebibliography}{99}

\bibitem[Bridgeford et al.(2026)]{bridgeford2026}
Bridgeford EW, Campbell ID, Chen Z, Lin Z, Ritz H, Vandekerckhove J, Poldrack RA.
Twelve quick tips for AI-assisted coding in science.
\textit{PLOS Computational Biology} 2026;22(7):e1014428.
\url{https://doi.org/10.1371/journal.pcbi.1014428}

\bibitem[Sandve et al.(2013)]{sandve2013}
Sandve GK, Nekrutenko A, Taylor J, Hovig E.
Ten simple rules for reproducible computational research.
\textit{PLOS Computational Biology} 2013;9(10):e1003285.

\bibitem[Wilson et al.(2014)]{wilson2014}
Wilson G, Aruliah DA, Brown CT, et al.
Best practices for scientific computing.
\textit{PLOS Biology} 2014;12(1):e1001745.

\bibitem[Poldrack(2024)]{poldrack2024}
Poldrack RA. Better code, better science. 2024.
\url{https://poldrack.github.io/BetterCodeBetterScience/}

\bibitem[van Burgsteden et al.(under review)]{vanburgsteden-thresholds}
van Burgsteden M, Geurs KT, Grigolon AB, Ulak MB.
Either you can reach it or you cannot: threshold microfoundations for
accessibility, applied to shared bicycle provision in Utrecht.
Under review.

\end{thebibliography}

\end{document}
